% EDITORIAL QUESTION: How does an external magnetic field drive topological gap opening and first-order phase transitions, and what laws govern the response?
\section{Zeeman Field Response, Cubic Dirac Mass, and First-Order Multistability}
\label{sec:field_response}

\subsection{First-Order Quantum Phase Transition and Hysteresis}
We apply an external magnetic field along the high-symmetry threefold axis:
\begin{equation}
\label{eq:field_111}
\mathbf{h} = \frac{h}{\sqrt{3}}(1, 1, 1)^T.
\end{equation}
As $h$ increases, the system undergoes a sharp quantum phase transition from the gapless/chiral spin liquid phase to the trivial field-polarized phase:

\begin{proposition}[First-Order Critical Boundary and Hysteresis]
\label{prop:first_order_transition}
Under numerical forward continuation ($h \nearrow$) and backward continuation ($h \searrow$) along the [111] field axis, the self-consistent Hartree-Fock solutions exhibit a pronounced first-order hysteresis loop:
\begin{enumerate}
\item \textbf{Thermodynamic Transition:} The ground-state energies cross at the critical field:
\begin{equation}
\label{eq:critical_field}
h_c^* / |K| = 0.3325 \pm 0.0005.
\end{equation}
\item \textbf{Metastable Coexistence Window and Domain-Wall Tension:} The QSL state persists metastably up to the forward spinodal breakdown $h_{c2}^+ = 0.400 \pm 0.005 |K|$, while the polarized state survives as a local minimum down to $h_{c1}^- = 0.0$. Near coexistence ($h = 0.335$), a planar QSL--polarized interface exhibits a strictly positive domain-wall surface tension $\sigma_{\mathrm{DW}} = +0.042 |K| > 0$ with interface healing width $\xi_{\mathrm{DW}} \approx 1.5 a$.
\item \textbf{Magnetization Discontinuity:} At $h_c^*$, the induced magnetization jumps discontinuously:
\begin{equation}
\Delta M \equiv M_{\mathrm{pol}}(h_c^*) - M_{\mathrm{QSL}}(h_c^*) = 0.412 \pm 0.005,
\end{equation}
confirming a discontinuous first-order transition with energy hysteresis barrier $\Delta E = 0.057 |K|$.
\item \textbf{Post-Spinodal KKT Saddle Audit:} Immediately above the upper spinodal ($h \in [0.400, 0.454]$), unstabilized 60-step continuation experiences transient excursions into unphysical gauge-violating saddle regions (\texttt{conv = False}, $\epsilon_{\mathrm{gauge}} \in [3.3 \times 10^{-4}, 3.58 \times 10^{-1}]$, producing spurious downward energy spikes to $E_{\mathrm{site}} = -1.1904$ and Bloch-ball violations $m_{111} = 1.1249 > 1$ at $h = 0.450$). Enforcing strict KKT convergence ($\epsilon_{\mathrm{gauge}} < 10^{-9}$) and Bloch-ball admissibility eliminates these artifacts, collapsing the state at $h_{c2}^+ = 0.400$ directly onto the physical polarized branch ($E_{\mathrm{site}}(h) = -h - 0.5000$, $m_{111} = 1.0000$, $\epsilon_{\mathrm{gauge}} \le 5.6 \times 10^{-15}$).
\end{enumerate}
\end{proposition}

\subsection{Cubic Majorana Dirac Mass and Chiral Perturbation Theory}
In the weak-field regime ($h < h_c^*$), the external field breaks time-reversal invariance $\mathcal{T}$, lifting the gapless Dirac nodes at the Brillouin zone corners $\mathbf{K}, \mathbf{K}'$:

\begin{proposition}[Symmetry Structure and Empirical Cubic Scaling Heuristic (\texttt{IDEA-001} Open)]
\label{thm:cubic_gap_law}
Let $\Delta_K(\mathbf{h}) \equiv \varepsilon_5(\mathbf{K}) - \varepsilon_4(\mathbf{K})$ denote the direct single-particle excitation gap evaluated at $\mathbf{K}$ in self-consistent 4-flavor Majorana mean-field theory. Under discrete $C_3$ bond permutation and time-reversal $\mathcal{T}$ ($\mathbf{h} \mapsto -\mathbf{h}$, $\mathbf{K} \leftrightarrow -\mathbf{K}$):
\begin{enumerate}
\item Quadratic polynomials $h_\alpha h_\beta$ are $\mathcal{T}$-even and shift the Dirac cone centers without generating a $\mathcal{T}$-breaking mass gap at linear level in a minimal two-band effective model. While $\mathcal{T}$-oddness and $C_3$ invariance alone do not uniquely rule out $h_x + h_y + h_z$ (which requires an explicit representation-level analysis of the sublattice/matter-gauge hybridization), numerical evaluation in the isotropic $[111]$ field window $h \in [0.005, 0.150]$ on a $24\times 24$ mesh shows dominant $\mathcal{O}(h^3)$ gap opening driven by the chiral triple product $\mathcal{I}^{(3)}(\mathbf{h}) = h_x h_y h_z = h^3/(3\sqrt{3})$.
\item Establishing an exact closed-form analytical series $\Delta_K(h) = c_{\mathrm{MF}}(h/|K|)^3 + \mathcal{O}(h^5)$ from the coupled HFB-KKT Bloch generator remains open under \texttt{IDEA-001}.
\end{enumerate}
\end{proposition}

\begin{remark}[Operator-Level Perturbation Mapping and Mean-Field Reconciliation]
In third-order Rayleigh-Schr\"odinger / Feshbach perturbation theory of the spin-1/2 Kitaev model \cite{Kitaev2006}, an external Zeeman field induces an effective three-spin chiral interaction on each triad of bonds meeting at vertex $k$:
\begin{equation}
\label{eq:kitaev_kappa_exact}
H^{(3)} = -\kappa \sum_{\langle\langle j, k, l \rangle\rangle} \sigma_j^\alpha \sigma_k^\gamma \sigma_l^\beta.
\end{equation}
Projecting into the unperturbed flux-free physical subspace $P_{\mathrm{ff}}$, the local gauge constraint ($\mathcal{D}_k = +c_k^0 c_k^x c_k^y c_k^z = +1$ in Eqs.~\eqref{eq:spin_algebra}--\eqref{eq:parity_projector}, identified with Kitaev's $D_k^{\mathrm{Kitaev}} = b_k^x b_k^y b_k^z c_k = +1$ via the odd-parity map $(c_k^0, c_k^1, c_k^2, c_k^3) = (-c_k, b_k^x, b_k^y, b_k^z)$ so that $\sigma_k^\alpha P_{\mathrm{ff}} = +\ii c_k^0 c_k^\alpha P_{\mathrm{ff}} = +\ii b_k^\alpha c_k$) contracts the internal gauge Majoranas into background $\mathbb{Z}_2$ links, yielding the leading operator channel structure:
\begin{equation}
\label{eq:three_spin_majorana_projection}
P_{\mathrm{ff}} \left(\sigma_j^\alpha \sigma_k^\gamma \sigma_l^\beta\right) P_{\mathrm{ff}} = \epsilon^{\alpha\beta\gamma} (\ii c_j c_l).
\end{equation}
Thus, the three-spin interaction maps to next-nearest-neighbor matter Majorana hopping, opening a canonical Dirac gap $\Delta_{\mathrm{canonical}} = 12\sqrt{3}\kappa$. Two important analytical caveats (\texttt{K-019-A} and \texttt{IDEA-001}) govern any quantitative comparison:
(i) In Kitaev's static single-mode approximation \cite{Kitaev2006} (replacing the full two-vison plus matter Feshbach--Schur resolvent $Q(E_0 - H_K)^{-1}Q$ by the isolated two-vison threshold $\Delta_{\mathrm{flux}} \approx 0.264 |K|$), $\kappa_{\mathrm{static}} \approx h_x h_y h_z / \Delta_{\mathrm{flux}}^2$ gives only the illustrative static estimate $\Delta_{\mathrm{static}} = \frac{4}{\Delta_{\mathrm{flux}}^2} h^3 \approx 57.40 h^3/|K|^2$. Moreover, vanishing first-order diagonal projection ($PVP = 0$) kills $H_{\mathrm{eff}}^{(1)}$ but does \emph{not} imply $H_{\mathrm{eff}}^{(2)} = PV Q(E_0 - Q H_0 Q)^{-1} Q VP = 0$; rigorous exact evaluation requires integrating the full-$Q$ resolvent across physical-parity eigenstates (tracked independently via \texttt{kitaev-walk-geometry} and repository Issues \texttt{\#10}, \texttt{\#12}).

In our self-consistent 4-flavor Majorana mean-field theory, finite-mesh numerical evaluation across $h \in [0.005, 0.150]$ yields the empirical power-law fit:
\begin{equation}
\label{eq:mf_gap_numerical_fit}
\Delta_K^{\mathrm{MF}}(h) = (2.5921 \pm 0.0034) \cdot \left(\frac{h}{|K|}\right)^{2.9936 \pm 0.0006},
\end{equation}
whose closed-form analytical coefficient derivation remains an open question (\texttt{IDEA-001}). Within Gaussian \texttt{scMFT}, the structural relationship between Eq.~\eqref{eq:three_spin_majorana_projection} and the fitted scale $\Delta_K^{\mathrm{MF}}$ is governed by three operator mechanisms:
\begin{enumerate}
\item \textbf{Bloch Generator Normalization (Factor of 4):} In the single-particle quadratic representation $H = \frac{\ii}{4}\mathbf{c}^T A \mathbf{c}$, the Bloch Hamiltonian is defined as $H_{\mathrm{Bloch}}(\mathbf{k}) = \frac{\ii}{4} A(\mathbf{k})$, scaling reported eigenvalues by $1/4$ relative to the canonical physical generator $\ii A$. Restoring canonical mode normalization yields the fitted scale $\Delta_{\mathrm{canonical}}^{\mathrm{MF}}(h) = 4 \Delta_K^{\mathrm{MF}}(h) \approx 10.3684 h^3/|K|^2$.
\item \textbf{Two-Vison Continuum vs. Flat Gauge-Band Denominators:} Whereas a physical Pauli operator $\sigma_i^\gamma = \ii c_i^0 c_i^\gamma$ flips a $\mathbb{Z}_2$ bond link and excites the non-local two-vison plus matter continuum above $\Delta_{\mathrm{flux}} \approx 0.264 |K|$, unprojected Hartree-Fock factorization preserves uniform links and excites localized flat-band gauge Majoranas $c_i^\alpha$ with two-particle gap $2\Delta_{\mathrm{gauge}} = 2|K|\bar{U} \approx 1.04973 |K|$ (giving a bare squared scale ratio $(2\Delta_{\mathrm{gauge}}/\Delta_{\mathrm{flux}})^2 \approx 15.8106$).
\item \textbf{Self-Consistent KKT Multiplier and Fock Vertex Dressing:} In 4-flavor HFB theory, the bare Zeeman matrix $-\sum_\gamma h_\gamma N^\gamma$ only couples $c^0 \leftrightarrow c^\gamma$ and cannot by itself mix distinct gauge flavors $c^\alpha \leftrightarrow c^\beta$ at the intermediate vertex. Instead, the $\mathcal{O}(h^3)$ chiral Dirac mass is generated by the self-consistent $\mathcal{O}(h)$ response: the linear KKT Lagrange multiplier $\lambda_s^\gamma G^\gamma$ ($G^\gamma_{\alpha\beta} = -\epsilon_{\gamma\alpha\beta}$, mixing $c^\alpha \leftrightarrow c^\beta$) together with the off-diagonal Fock bond self-energy $\delta U_b = \mathcal{O}(h)$.
\end{enumerate}
\end{remark}

\subsection{Lagrange Multiplier Gauge Response Theory}
Because an external Zeeman field breaks flat-band gauge neutrality, non-zero Lagrange multipliers $\boldsymbol{\lambda}_s^*$ are generated to enforce physical state representability $\langle \mathbf{G}_s \rangle = \mathbf{0}$:

\begin{proposition}[Gauge Susceptibility and Linear Counter-Torque Response]
\label{prop:lagrange_law}
Applying the Implicit Function Theorem to the stationarity condition $\mathbf{g}_s(\boldsymbol{\lambda}^*(\mathbf{h}), \mathbf{h}) \equiv \mathbf{0}$ yields the exact linear response derivative:
\begin{equation}
\label{eq:lagrange_multiplier_linear_response}
\left. \frac{\mathrm{d}\boldsymbol{\lambda}_s^*}{\mathrm{d}\mathbf{h}} \right|_{\mathbf{h}=\mathbf{0}} = -\boldsymbol{\chi}_{\mathrm{gauge}}^{-1} \boldsymbol{\chi}_{G, S},
\end{equation}
where $\boldsymbol{\chi}_{\mathrm{gauge}} \equiv \frac{\partial \langle \mathbf{g}_s \rangle}{\partial \boldsymbol{\lambda}_s}$ is the local gauge susceptibility matrix and $\boldsymbol{\chi}_{G, S} \equiv \frac{\partial \langle \mathbf{g}_s \rangle}{\partial \mathbf{h}}$ is the cross-susceptibility describing Zeeman-induced gauge polarization.
In the weak-field regime along the symmetric axis $\mathbf{h} = \frac{h}{\sqrt{3}}(1, 1, 1)^T$, finite-mesh self-consistent solutions obey the leading-order empirical regression (with closed-form analytical prefactor left open in \texttt{IDEA-001}):
\begin{equation}
\label{eq:lagrange_multiplier_fit}
\|\boldsymbol{\lambda}^*\|_2 \approx \sqrt{5} \cdot h \approx 2.236 \cdot h,
\end{equation}
with relative residuals bounded by $< 0.8\%$ across $h \le 0.10$.
\end{proposition}
This establishes that local gauge projection acts as an effective counter-torque whose restoring strength grows proportionally to the external symmetry-breaking field.

\subsection{Hellmann-Feynman Thermodynamic Consistency}
To verify the variational integrity of the solutions, we test the Hellmann-Feynman theorem relating magnetization to energy derivatives:
\begin{equation}
\label{eq:hellmann_feynman}
M_{\mathrm{HF}}(h) = -\frac{\partial E_{\mathrm{MF}}}{\partial h}.
\end{equation}
Across all fields $h \in [0, 1.0]$, central finite differences $(\Delta h = 10^{-4})$ agree with quantum expectation values $M_{\mathrm{obs}} = \frac{1}{N_{\mathrm{sub}}} \sum_s \langle \sigma_s^{[111]} \rangle$ to within $2.1 \times 10^{-6}$, confirming exact variational stationarity.
